\documentclass[10pt,a4paper]{amsart}
\usepackage[margin=19mm]{geometry}
\usepackage{amsmath,amssymb,mathtools}
\usepackage[scr=rsfs,cal=euler]{mathalfa}
\usepackage{xfrac}
\usepackage[unicode,hidelinks,bookmarks=false]{hyperref}
\newcommand{\ie}{i.e.\ }
\newcommand{\cf}{cf.\ }
\newcommand{\resp}{respectively\ }
\newcommand{\Subsection}{\subsection}
\providecommand{\nobreakdash}{\nobreak\hbox{-}}
\setlength{\emergencystretch}{2em}
\multlinegap0pt
\usepackage{amssymb}
\newtheorem{lemma}{Lemma}
\newcommand{\Var}{\operatorname{Var}}
\title{Variance of random greedy independent sets in triangle-free graphs}
\author{Mubin Shaikh}
\date{Source: arXiv:2609.14826v1 (CC BY 4.0)}
\begin{document}
\maketitle

\noindent\emph{Source and licence.} Variance of random greedy independent sets in triangle-free graphs. Version arXiv:2609.14826v1 (CC BY 4.0), distributed by arXiv under the Creative Commons Attribution 4.0 International licence, http://creativecommons.org/licenses/by/4.0/, as recorded in the arXiv licence metadata for this exact version and shown on the abstract page https://arxiv.org/abs/2609.14826v1. Full licence notice: \texttt{licenses/arxiv-2609.14826-CC-BY-4.0.txt}.\par\smallskip

\noindent\textit{Editorial scope.} Counted unit: the article's lemma lem:centered (source0.tex lines 113--123). The article's complete original proof of it is delivered with it (source0.tex lines 125--162, one \texttt{proof} environment of 38 lines). No mathematics is written, repaired or re-derived: the statement and the proof are the article's own lines, in the article's own notation, and no retained line is substituted or re-flowed.  Retained uncounted as the source-local context the proof needs: lines 77--79; lines 88--90; lines 107--110.  Nothing was omitted from any retained span, so the package carries no omission record.  No retained line contains a citation, so the wrapper carries no bibliography and no bibliography entry.  No retained line contains a figure environment, an image-inclusion command or drawing code, so no image is delivered and the package declares no asset dependency.  Editorial formatting only, in addition: an AMS \texttt{amsart} preamble that restates the article's own declarations and macros used by the retained lines, namely the environments \texttt{lemma} on separate counters, together with the standard packages those lines need.  The retained environments print in this document as Lemma 1 in order of appearance, whereas the article prints the same items with the numbering of its own full document.  The complete native source of the article as distributed in the arXiv e-print for this exact version is bundled unchanged as \texttt{sources/210362/source0.tex}.
\begin{equation}\label{eq:mean}
 \mu(G)=\frac1n\sum_u a_u(G),\qquad \sum_u b_u(G)=0.
\end{equation}

\begin{equation}\label{eq:lipschitz}
 |\mu(G[A])-\mu(G[B])|\le |A\triangle B|.
\end{equation}

\begin{equation}\label{eq:crude-bias}
 |b_x(G)|=|\mu(G)-\mu(G-N_G(x))|\le d_G(x),
 \qquad \mu(G-N_G(x))=1+\mu(R_x).
\end{equation}

\begin{lemma}[Centered first-choice estimate]\label{lem:centered}
For every triangle-free graph $G$ of order $n\ge2$ and every $x\in V(G)$,
\begin{equation}\label{eq:centered}
 |b_x(G)|\le c_n d_G(x).
\end{equation}
For a singleton, $b_x(G)=0$. Hence, for a uniform vertex $u$,
\begin{equation}\label{eq:between}
 \Var_u(a_u(G))=\frac1n\sum_u b_u(G)^2
 \le\frac{c_n^2}{n}\sum_u d_G(u)^2.
\end{equation}
\end{lemma}

\begin{proof}
The singleton case is immediate; assume $n\ge2$.

Fix $x$, put $d=d_G(x)$, and let $U=V(G)\setminus N_G[x]$. We first establish
\begin{equation}\label{eq:recurrence}
 n b_x(G)=\sum_{u\in N_G(x)}\bigl(\mu(R_u)-\mu(R_x)\bigr)
                  +\sum_{u\in U}b_x(R_u).
\end{equation}
For $u\in U$, deleting the two closed neighborhoods commutes:
\[
 R_u-N_{R_u}[x]=G-(N_G[u]\cup N_G[x])=R_x-N_{R_x}[u].
\]
Thus $\mu(R_u)=1+\mu(R_x-N_{R_x}[u])+b_x(R_u)$. Sum over $u\in U=V(R_x)$ and apply \eqref{eq:mean} to $R_x$ to obtain
\[
 \sum_{u\in U}\mu(R_u)=|U|\mu(R_x)+\sum_{u\in U}b_x(R_u).
\]
When $U=\varnothing$ both sides are zero. Substitution into $n\mu(G)=n+\sum_u\mu(R_u)$ gives \eqref{eq:recurrence}.

If $u$ is adjacent to $x$, triangle-freeness gives $N_G[u]\cap N_G[x]=\{u,x\}$. By \eqref{eq:lipschitz},
\[
 |\mu(R_u)-\mu(R_x)|\le d_G(u)+d-2.
\]
For $u\in U$, apply \eqref{eq:crude-bias} to $R_u$ to get $|b_x(R_u)|\le d_{R_u}(x)$. A count of ordered pairs $(u,w)$, where $w\in N_G(x)$ is removed by deleting $N_G[u]$, gives
\begin{equation}\label{eq:degreecount}
 \sum_{u\in U}d_{R_u}(x)
 =(n-d-1)d-\sum_{w\in N_G(x)}(d_G(w)-1).
\end{equation}
Indeed, for each such $w$, every vertex of $N_G(w)\setminus\{x\}$ lies in $U$, because there are no triangles. The count allows different common neighbors, so it does not exclude four-cycles.

Taking absolute values in \eqref{eq:recurrence} now gives
\begin{align*}
 n|b_x(G)|
 &\le \sum_{u\in N_G(x)}(d_G(u)+d-2)
 +(n-d-1)d-\sum_{w\in N_G(x)}(d_G(w)-1)\\
 &=d(n-2).
\end{align*}
This proves \eqref{eq:centered}, also for $n=2$ and for isolated $x$. Equation \eqref{eq:between} follows from \eqref{eq:mean}.
\end{proof}

\end{document}
